% !TeX root = main.tex
%
% Related work for the Pavillon 3DGS report. Organised by the axes on which our
% capture is hard, so each family is judged by whether it addresses OUR failure
% modes rather than by general popularity.

\section{Related Work}
\label{sec:related}

Our capture is adversarial along four specific axes: it is \emph{single-sided}
(the camera never orbits the object), \emph{low-parallax}, \emph{close-range} on a
\emph{near-planar relief}, and it is assembled from handheld clips with
independent auto-exposure. We therefore organise the literature by those axes and,
where we tested a family, report what it did on this data.

%---------------------------------------------------------------------------
\subsection{Primitive-based radiance fields and frameworks}

3DGS \cite{kerbl2023gaussian} replaced NeRF's ray-marched implicit field
\cite{mildenhall2020nerf} with explicit anisotropic primitives rasterised by EWA
splatting \cite{zwicker2001ewa,zwicker2001surface}, trading a continuous
representation for real-time rendering and a far cheaper optimisation. Two open
implementations now dominate practice: \texttt{gsplat} \cite{ye2024gsplat}, a
CUDA-backed library exposing the rasteriser and density-control strategies, and
Nerfstudio \cite{tancik2023nerfstudio}, whose \emph{Splatfacto} model wraps
\texttt{gsplat} with a configurable training pipeline. We build directly on the
packaged \texttt{gsplat} APIs rather than on Nerfstudio: the pipeline needs
per-stage checkpointing, preemption handling and fingerprint-driven reruns on a
Slurm cluster, which is orchestration Nerfstudio does not target, and avoiding the
Nerfstudio dependency also avoids \texttt{tiny-cuda-nn}, a known pain point on
Blackwell (\texttt{sm\_120}).

Anti-aliasing is orthogonal but relevant to close-range capture: Mip-Splatting
\cite{yu2024mipsplatting} adds 3D smoothing and a 2D Mip filter so primitives
observed at sampling rates far from training do not produce erosion or dilation
artefacts. Our capture varies sharply in scale between contextual and extreme
close-up views, so we enable \texttt{gsplat}'s antialiased rasterisation mode
throughout.

%---------------------------------------------------------------------------
\subsection{Density control: from heuristics to sampling}
\label{sec:rw-density}

Vanilla 3DGS grows the model with a hand-tuned adaptive density control (ADC):
clone or split wherever the view-space positional gradient exceeds a threshold,
periodically reset opacities, prune the transparent. The thresholds are heuristic
and, as we found empirically (\S\ref{sec:theory} and our capacity experiments),
their interaction with scene scale is fragile.

Two lines of work replace or restructure this. \textbf{Structured
representations} anchor primitives to a coarser scaffold: Scaffold-GS
\cite{lu2024scaffoldgs} distributes local Gaussians around anchor points and
predicts their attributes on the fly from view direction and distance, which
suppresses the redundant floaters that unconstrained ADC produces; Octree-GS
\cite{ren2024octreegs} extends this with an LOD hierarchy so the level of detail
matches the observation scale.

\textbf{3DGS-MCMC} \cite{kheradmand2024mcmc} is the more fundamental reframing and
the most relevant to our findings. It reinterprets the Gaussians as samples drawn
from a distribution over scenes, turns the optimiser into Stochastic Gradient
Langevin Dynamics by injecting noise, and rewrites cloning and pruning as
\emph{state transitions that approximately preserve sample probability} — replacing
``dead'' primitives by relocating them onto live ones rather than heuristically
splitting. Crucially, it reports better reconstructions \emph{at an identical
Gaussian budget}, and it makes the budget an explicit parameter rather than an
emergent consequence of thresholds.

This connects directly to our own result. We measured a monotone preference for
\emph{fewer} primitives on this capture: reducing \texttt{cap\_max} from 1.5\,M to
375\,k raised test PSNR by 1.8\,dB, improved SSIM, and eliminated the
catastrophic-view tail, while doubling it to 3.0\,M degraded every metric. The
compaction literature — Compact-3DGS \cite{lee2024compact3dgs}, GaussianSpa
\cite{zhang2024gaussianspa} — reports that aggressive pruning preserves or improves
quality while shrinking models by large factors, which is consistent, though that
work is usually motivated by memory rather than generalisation. The sparse-view
literature makes the generalisation argument explicitly: DropGaussian
\cite{park2025dropgaussian} randomly removes primitives during training as a
\emph{structural regulariser} for few-view inputs, and recent analysis quantifies
\emph{co-adaptation} between primitives as a distinct sparse-view failure mode
\cite{coadaptation2025}. Our capacity sweep is an independent, purely empirical
instance of the same phenomenon, obtained without any architectural change: on a
low-parallax capture, primitive count behaves as a regularisation hyper-parameter,
not a quality dial. Adopting 3DGS-MCMC would be the principled next step, since it
makes exactly this budget explicit.

%---------------------------------------------------------------------------
\subsection{Sparse and few-view reconstruction}

When views are few or share little parallax, the photometric objective is
underdetermined: many primitive configurations reproduce the training images while
placing mass at wrong depths. The dominant remedy is to import geometry from a
monocular depth network \cite{ranftl2022midas,yang2024depthanythingv2}, used
scale-invariantly because such predictions are only defined up to an affine
transform \cite{eigen2014depth}. FSGS \cite{zhu2024fsgs} densifies guided by depth
and proximity; SparseGS \cite{xiong2024sparsegs} combines depth with a
diffusion-based score to remove background collapse; DNGaussian
\cite{li2024dngaussian} contributes global-local depth normalisation so that
absolute depth error does not dominate local shape. CoR-GS \cite{zhang2024corgs}
takes a different route, training two models and using their disagreement to
identify unreliable geometry.

We implement the depth-prior family (DepthAnything-v2 with a Pearson-correlation
loss, \S\ref{sec:sparse}) and measure it honestly: on this capture the prior is
\emph{nearly free} in PSNR terms (an ablation with it disabled recovers only
0.26\,dB) and slightly \emph{increases} floater removal — but it is not what
determines final quality. Resolution and capacity dominate it by a wide margin.
This is a useful negative-ish result: depth priors are routinely presented as the
answer to sparse-view degradation, yet on a capture whose views are numerous but
badly distributed, the binding constraints were elsewhere.

A newer direction sidesteps SfM entirely, initialising from learned stereo:
DUSt3R \cite{wang2024dust3r} and VGGT \cite{wang2025vggt} regress dense geometry
and poses directly, and InstantSplat \cite{fan2024instantsplat} couples such an
initialisation to 3DGS for very sparse inputs. For our data the SfM stage was not
the bottleneck once global mapping was used (below), so we did not pursue it;
it remains the natural fallback for captures where COLMAP/GLOMAP fail outright.

%---------------------------------------------------------------------------
\subsection{Camera calibration and pose refinement}

Incremental SfM \cite{schoenberger2016sfm} registers images greedily and can
abandon weakly connected regions — exactly the regime of a low-overlap capture.
Global SfM re-estimates all poses jointly; GLOMAP \cite{pan2024glomap} performs
global rotation averaging followed by joint global positioning. Substituting
GLOMAP for incremental COLMAP on identical features and matches raised
registration from 82 to 181 of 193 images on our data, a $2.2\times$ coverage gain
and the single largest early improvement we measured.

Where poses remain inaccurate, they can be refined jointly with the scene, as in
BARF \cite{lin2021barf}, whose central observation is that such optimisation must
be coarse-to-fine to avoid poor local minima; COLMAP-Free 3DGS
\cite{fu2024colmapfree} pushes this to joint pose-and-scene estimation without SfM
at all. We implemented per-camera SE(3) refinement and found it \emph{harmful}
here ($-1$\,dB): GLOMAP's poses were already sub-pixel (0.92\,px mean reprojection
error), so there was no calibration error to recover, and refining only the
training cameras — held-out poses must not be optimised, or the metric is
contaminated — introduced a small relative misalignment between the reconstructed
scene and the un-refined evaluation cameras. Pose refinement is a remedy for bad
poses, and ours were not bad.

%---------------------------------------------------------------------------
\subsection{Surfaces, meshes and hybrid representations}

A radiance field is not a surface, and our subject is a carved relief whose value
is geometric. Several families address this. \textbf{Flattening the primitive}:
2DGS \cite{huang20242dgs} replaces ellipsoids with oriented disks that are surface
elements by construction, and SuGaR \cite{guedon2024sugar} regularises 3D Gaussians
towards a surface before Poisson meshing. \textbf{Reformulating the depth or
opacity}: Gaussian Opacity Fields \cite{yu2024gof} extracts level sets of an
opacity field with Marching Tetrahedra, avoiding Poisson artefacts in unbounded
scenes; PGSR \cite{chen2024pgsr} compresses Gaussians to planes with multi-view
geometric consistency; RaDe-GS \cite{zhang2024radegs} derives a rasterised depth
that is unbiased for the level set. \textbf{Hybrid supervision}: GSDF
\cite{yu2024gsdf} runs a 3DGS branch and a neural SDF branch that supervise each
other — the splatting branch's rendered depth guides SDF ray sampling, and the SDF
in turn guides density control to grow primitives near the surface and prune those
that drift off it. GSDF is the most attractive of these for our problem, because
its density control is driven by a \emph{geometric} signal rather than a
photometric gradient, which is precisely the signal our capture lacks.

We evaluated 2DGS and report a clear negative result: it collapses to flat,
featureless renders (PSNR $\approx$13) on this capture. The explanation is
structural rather than incidental — a disk must recover an orientation as well as a
position, and normals are constrained mainly by observing a surface from multiple
directions, which a single-sided capture does not supply. The added constraint
removes degrees of freedom the photometric loss was using without supplying the
multi-view evidence needed to fix the rest. We therefore extract geometry by
TSDF-fusing \cite{curless1996tsdf} the depth composited by the \emph{3D} Gaussian
model, accepting the known bias of volumetric primitives, and obtain a usable mesh
of the panel. GSDF or GOF would be the principled improvements.

%---------------------------------------------------------------------------
\subsection{Appearance variation across captures}

Merging clips is the cheapest way to buy coverage, but each handheld clip carries
its own auto-exposure and white-balance response, and 3DGS has no parameter
expressing ``same radiance, different camera response''. NeRF-W
\cite{martinbrualla2021nerfw} introduced per-image appearance latents optimised
jointly with the scene, in the GLO style \cite{bojanowski2018glo}. The idea has
since been ported to splatting: Splatfacto-W \cite{xu2024splatfactow} adds
per-Gaussian neural colour features, per-image appearance embeddings and a
spherical-harmonic background to Nerfstudio's Splatfacto; WildGaussians
\cite{kulhanek2024wildgaussians} combines appearance modelling with DINO features
for occlusion robustness. A lighter-weight alternative now standard in Nerfstudio
is the bilateral grid of \citet{wang2024bilateral}, which learns a smooth
low-resolution 3D colour-transform grid per image — more expressive than a global
affine map while still spatially constrained.

We implement the low-capacity end of this spectrum deliberately: a per-image latent
decoded to a \emph{global} affine colour transform. The capacity argument matters
more than the architecture. A per-pixel correction could memorise each training
view and destroy the multi-view constraint; our transform applies identically at
every pixel, so it commutes with any spatial permutation of the image and therefore
carries exactly zero spatial information. It can absorb exposure, gain and
white-balance drift and provably nothing else. A bilateral grid sits between these
extremes and would be the natural upgrade if per-clip vignetting or spatially
varying response mattered.

Our merge nevertheless underperformed the best single-clip model. Diagnosing it
required first fixing a measurement bug — evaluation must apply the appearance
correction, and must do so using a latent that does not depend on the held-out
image, or the metric is either penalised or contaminated. With per-clip latents the
merged model reaches 22.0\,dB against 24.9\,dB single-clip. The residual gap is not
photometric: the second clip enlarges the reconstructed volume while contributing
few views, and the capacity effect of \S\ref{sec:rw-density} then dominates.

%---------------------------------------------------------------------------
\subsection{Positioning}

Relative to this literature our contribution is not a new representation but a
measured account of \emph{which} interventions matter for one well-characterised
hard capture, with negative results reported alongside positive ones. Global SfM
and resolution dominated; capacity reduction was the largest late gain; depth
priors were nearly free but not decisive; 2DGS, multi-clip merging and pose
refinement each degraded quality for reasons we can attribute. The methods we did
not implement — 3DGS-MCMC for principled budget control, GSDF or GOF for geometry,
bilateral grids for appearance — are exactly the ones this evidence points to.
